\begin{abstract}

Low-cost programmable mixed-signal platforms provide an attractive foundation for scientific instrumentation, but the limited precision of their configurable analog building blocks often constrains analog front-end performance. This paper proposes a Digitally Assisted Analog Front-End architecture that temporarily exploits the embedded mixed-signal resources of a programmable system-on-chip (PSoC) to improve the operating conditions of the analog signal chain while preserving a conventional acquisition path during normal operation.

As a proof of architectural feasibility, the proposed architecture is demonstrated through an automatic foreground DC-offset calibration that sequentially recenters the programmable-gain amplifier, band-pass filter, summing stage, and low-pass filter using the on-chip ADC, DAC, and digital filtering resources.

The architecture was experimentally evaluated on two independently assembled geophone acquisition boards under eight board--gain configurations. Relative to a common fixed-reference baseline, the calibrated low-pass output reduced the residual DC offset from 2.9--45.8\% to 0.4--3.9\% of the nominal ADC input span using a single controller configuration across all reported conditions.

These results demonstrate the feasibility of the proposed Digitally Assisted Analog Front-End architecture as a practical approach for improving the operating conditions of low-cost configurable analog hardware using embedded mixed-signal resources. Statistical validation over larger hardware populations and broader environmental characterization remain future work.

\end{abstract}

\begin{IEEEkeywords}
Digitally assisted analog design, analog front end, mixed-signal systems, geophone instrumentation, programmable system-on-chip (PSoC).
\end{IEEEkeywords}
